aset contains trajectories collected from a
partially trained policy, medium-replay contains the corresponding replay
buffer collected throughout training, and medium-expert combines medium
and expert behavior. Hopper-medium-replay is used for the primary mechanism
diagnostic because its broader behavioral coverage provides a wider range of
learner-relative distances, while all three Hopper datasets are retained for
the task-level performance evaluation.

\subsection{Countdown Task and Puzzle Dataset}
\label{app:countdown_dataset}

Countdown is an arithmetic generation task in which the policy receives
a multiset of integers and a target value, and must generate an
arithmetic expression that evaluates exactly to the target. A valid
output must use every supplied number exactly once and may contain only
addition, subtraction, multiplication, division, and parentheses.

Let
$x=(\mathcal N,T)$
denote a puzzle, where
$\mathcal N$
is the supplied multiset of numbers and
$T$
is the target. The policy action is a complete token sequence
$y$ representing the generated expression. Task success is determined
by a deterministic verifier:
\begin{equation}
    R(x,y)
    =
    \mathbf 1
    \left[
        y
        \text{ is syntactically valid, uses }
        \mathcal N
        \text{ exactly once, and evaluates to }T
    \right].
    \label{eq:countdown_reward}
\end{equation}
Because multiple expressions can solve the same puzzle, evaluation is
based on verifier success rather than exact string matching.

The puzzle corpus is procedurally generated from valid arithmetic
expressions and divided into disjoint training, validation, and test
sets. Duplicate puzzles with the same number multiset and target are
removed across the corpus. Validation puzzles are used for checkpoint
selection, while the test split is reserved for final evaluation.

Countdown provides an external discrete-action setting with
autoregressive sequence generation and shared Transformer parameters.
For the diagnostic in Figure~\ref{fig:external_far_field}, unsuccessful
responses are verifier-matched and grouped by current mean-token surprisal.

\subsection{Training and Evaluation Protocols}

\section{Controlled Experimental Protocols}
\label{app:controlled_protocols}

This section specifies the interventions used in
Sections~\ref{sec:controlled_identification}
and~\ref{sec:repulsion_control}. The definitions of C-U1 and D-U1 are given in
Appendix~\ref{app:experimental_details}. Network architectures,
optimizer settings, training budgets, seeds, and evaluation frequencies
are reported in the protocol-specific subsections below.

\subsection{Shared Notation and Experimental Invariants}
\label{app:controlled_shared}

For a negative sample
$z=(s,a)$, define its raw repulsive update as
\begin{equation}
    g_\theta^-(z)
    =
    -\widehat A^-(z)
    \nabla_\theta
    \log\pi_\theta(a\mid s),
    \qquad
    \widehat A^-(z)
    =
    \max\{-\widehat A(z),0\}.
    \label{eq:controlled_negative_update}
\end{equation}
Learner-relative remoteness is measured by
\begin{equation}
    D_\theta(z)
    =
    -\log\pi_\theta(a\mid s),
    \label{eq:controlled_remoteness}
\end{equation}
or by its equivalent standardized Gaussian component when additive
density constants are irrelevant.

Unless explicitly stated otherwise, all branches of one controlled
comparison share:

\begin{enumerate}
    \item the same training and evaluation contexts;
    \item the same positive and negative samples;
    \item the same fixed rewards or advantage labels;
    \item the same initial policy parameters;
    \item the same minibatch order, optimizer, and training horizon;
    \item the same evaluation checkpoints and terminal classification
    rules.
\end{enumerate}

Only the variable named by the intervention is changed. Actions generated
under the same context are aggregated at the context level and are not
treated as independent contexts for statistical inference.


\subsection{Source-Isolation Protocol}
\label{app:source_isolation_protocol}

\paragraph{Purpose.}
The source-isolation experiment asks whether policy-relative remoteness
changes negative-update magnitude when sample quality is held fixed. It
identifies the source of a per-sample difference, not its causal effect on
training.

\paragraph{Continuous product-manifold construction.}
For each C-U1 context, we select multiple negative actions from a common
reward contour:
\begin{equation}
    R(s,a_i)
    =
    R(s,a_j),
    \qquad
    \widehat A(s,a_i)
    =
    \widehat A(s,a_j)
    <0,
    \label{eq:source_equal_quality}
\end{equation}
while ensuring that
\begin{equation}
    D_\theta(s,a_i)
    \neq
    D_\theta(s,a_j).
    \label{eq:source_different_remoteness}
\end{equation}
The reward coordinate and the policy-distance coordinate are therefore
independently controlled. This exact counterfactual comparison is the
reason for using a synthetic product-manifold construction: an external
offline dataset generally cannot provide two actions with precisely
matched task quality and independently prescribed distance from the
same current policy.

The comparison is performed both at policy initialization and after
positive-only pretraining. No negative sample used as a probe is applied
as an optimizer update during this protocol. Let \(u_\theta(s)\) denote
the policy-output coordinate used by the controlled probe, e.g.,
fixed-covariance Gaussian mean coordinates in C-U1 and direct logits in
D-U1. For every probe, we record
\begin{equation}
\begin{aligned}
I_{\mathrm{score}}(z)
&=
\left\|
\nabla_{u}\log \pi_{u}(a\mid s)
\right\|_2,\\
I_{\mathrm{sample}}(z)
&=
|\widehat A(z)|\, I_{\mathrm{score}}(z).
\end{aligned}
\label{eq:source_influence_metrics}
\end{equation}
Because \(|\widehat A|\) is identical across the matched actions,
differences in \(I_{\mathrm{sample}}\) arise from policy-score geometry
rather than advantage magnitude.

For neural implementations, we additionally record the full
parameter-gradient norm and the aggregate gradient obtained from samples
in the same remoteness range. These implementation-level actor-gradient
diagnostics are used as transfer checks, not as theorem-level score
measures.

We report near-to-far ratios of remoteness, score norm, single-sample
gradient norm, aggregate gradient norm, and gradient-direction
coherence. Ratios are first computed within each context and seed and
are then aggregated across the prespecified seeds.

\paragraph{Categorical counterpart.}
D-U1 uses unordered categorical actions with identical fixed negative
labels. The comparison varies the current probability assigned to the
negative action while preserving its reward role and advantage:
\begin{equation}
    \widehat A(s,i)
    =
    \widehat A(s,j)
    <0,
    \qquad
    \pi_\theta(i\mid s)
    \neq
    \pi_\theta(j\mid s).
    \label{eq:categorical_source_match}
\end{equation}
For a direct categorical logit parameterization, the score norm is
bounded and need not grow without limit as probability decreases.
The relevant far-field effect is instead persistence: the negative
update continues to increase the selected action's logit gap even after
its probability has become small. We therefore report probability,
surprisal, logit gap, score norm, and cumulative suppression over
repeated updates. This distinguishes categorical support suppression
from the unbounded Gaussian-distance amplification measured in C-U1.


\subsection{Causal-Transmission Protocol}
\label{app:causal_transmission_protocol}

\paragraph{Purpose.}
The source-isolation protocol identifies the source of abnormal update
magnitude but not its effect on task behavior. Causal transmission is tested in the nonlinear
Gaussian protocol, where the policy mean and variance evolve jointly,
and in the categorical support reconstruction of D-U1.

\paragraph{Dynamic near--far partition.}
At each intervention step, negative samples are partitioned using their
current-policy remoteness. Let
$\mathcal N_\theta$
and
$\mathcal F_\theta$
denote the prespecified near- and far-field sets. Their membership is
recomputed as the policy evolves rather than frozen at initialization.
The thresholds or quantiles are fixed before training and specified with the
corresponding environment.

\paragraph{Interventions.}
For a negative sample $z$, the following multipliers are applied to
$g_\theta^-(z)$:
\begin{equation}
\begin{aligned}
    w_{\mathrm{uncontrolled}}(z)
    &=1,\\
    w_{\mathrm{positive}}(z)
    &=0,\\
    w_{\mathrm{near\text{-}zero}}(z)
    &=
    \mathbf 1[z\notin\mathcal N_\theta],\\
    w_{\mathrm{far\text{-}zero}}(z)
    &=
    \mathbf 1[z\notin\mathcal F_\theta].
\end{aligned}
\label{eq:causal_zero_interventions}
\end{equation}

The far-cap intervention preserves the direction of a far-field update
but limits its detached magnitude to a reference level estimated from
near-field negative samples:
\begin{equation}
    w_{\mathrm{far\text{-}cap}}(z)
    =
    \begin{cases}
        \displaystyle
        \min\left\{
            1,
            \frac{C_{\mathrm{near}}}
                 {I_{\mathrm{sample}}(z)+\epsilon}
        \right\},
        & z\in\mathcal F_\theta,\\[3mm]
        1,
        & \text{otherwise},
    \end{cases}
    \label{eq:causal_far_cap}
\end{equation}
where $C_{\mathrm{near}}$ is the prespecified near-field update-magnitude
reference.

To distinguish selective intervention from a generic decrease in
negative optimization, the global-matched control applies one scalar
to every negative sample:
\begin{equation}
    w_{\mathrm{global}}(z)
    =
    \alpha_{\mathrm{match}}.
    \label{eq:causal_global_weight}
\end{equation}
The scalar is chosen using the training batch so that its detached aggregate
negative-update magnitude matches the corresponding selective
intervention:
\begin{equation}
    \alpha_{\mathrm{match}}
    \sum_{z\in\mathcal B^-}
        I_{\mathrm{sample}}(z)
    =
    \sum_{z\in\mathcal B^-}
        w_{\mathrm{selective}}(z)
        I_{\mathrm{sample}}(z).
    \label{eq:causal_budget_match}
\end{equation}
The matching coefficient is not differentiated through.

\paragraph{Causal contrasts.}
The primary causal contrast is between removing near-field and
far-field negative updates. If Near-zero leaves the dynamics unchanged
while Far-zero or Far-cap alters them, the effect cannot be attributed
to negative feedback in general. The global-matched branch quantifies
how much of the change follows from reducing total negative-update magnitude
and how much requires selecting samples by remoteness.

\paragraph{Temporal and terminal measurements.}
We record policy displacement, reward, negative-gradient magnitude,
mean dynamics and, for learnable-variance runs, scale dynamics, entropy or effective support, and the first
time at which each prespecified event occurs. The analysis checks whether
far-field update magnitude rises before policy drift and whether drift precedes
task-performance loss.

Terminal outcomes are classified separately as:

\begin{enumerate}
    \item \textbf{task-performance collapse}: a sustained loss of the
    prespecified task metric while optimization remains finite;
    \item \textbf{support or variance-boundary event}: categorical
    effective support or Gaussian variance crosses its prespecified
    boundary without requiring NaN or Inf;
    \item \textbf{numerical collapse}: a parameter, loss, gradient, or
    model output becomes NaN or Inf.
\end{enumerate}

A boundary event is not automatically counted as task-performance
collapse, and finite execution is not automatically counted as
convergence.

\subsection{Controlled Negative-Strength Sweep}
\label{app:negative_strength_sweep}

The controlled negative-strength sweep varies only the effective
negative-repulsion strength $q/p$, where $p$ denotes the aggregate
positive-attraction budget and $q$ denotes the aggregate
negative-repulsion budget. All sweep arms share the same data, fixed
labels, initialization, optimizer, training horizon, and event
thresholds. The positive-only baseline corresponds to $q/p=0$.

Held-out-context reward evaluates unseen-state generalization within the
C-U1 controlled protocol. Policy shift is the normalized displacement of
the learned policy from the positive-only target and visualizes how far
negative feedback pushes the policy beyond the imitation solution.
Task-performance collapse, support or variance-boundary events, and
NaN/Inf numerical failures are reported as separate event types.


\subsection{Controlled Taper Comparison}
\label{app:controlled_taper_protocol}

\paragraph{Controlled utility construction.}
In C-U1, the hidden optimum $a_\star(s)$ lies beyond the center of the
positive demonstrations $a_+(s)$. A local negative anchor is placed on
the opposite side of $a_+(s)$, so repulsion from this anchor points
toward $a_\star(s)$. Additional negative actions are placed on matched
reward contours at larger policy-relative distances. Increasing
negative strength therefore produces three observable regimes:
positive-only under-extrapolation, useful local extrapolation, and
far-field over-extrapolation or instability.

D-U1 implements the same logic with unordered semantic actions. The
hidden optimal target lies beyond the demonstrated positive direction,
while the local negative target lies in the opposite semantic
direction. Repelling the local negative action can therefore improve
the shared representation toward the hidden optimal action. Additional
remote negative actions create support pressure without changing their
fixed negative label. Evaluation on independently sampled contexts is
reported as held-out-context generalization.

Let
\[
x(D)
=
\left[
\frac{D-\tau}{c}
\right]_{+},
\qquad
r_x(D)=\sqrt{x(D)}.
\]
The Linear and Quadratic names refer to polynomial order in the
radialized coordinate \(r_\theta=\sqrt{x_\theta}\), rather than in the
squared-remoteness coordinate \(x_\theta\).
\paragraph{Weighting rules.}
For negative-sample remoteness $D$, the compared rules are
\begin{equation}
\begin{aligned}
    \omega_{\mathrm{positive}}(D)
    &=0,\\
    \omega_{\mathrm{uncontrolled}}(D)
    &=1,\\
    \omega_{\mathrm{global}}(D)
    &=\alpha,\\
    \omega_{\mathrm{hard}}(D)
    &=
    \mathbf 1[D\leq\tau],\\
    \omega_{\mathrm{Rec\text{-}L}}(D)
    &=
    \frac{1}{1+\lambda r_x(D)},\\
    \omega_{\mathrm{Rec\text{-}Q}}(D)
    &=
    \frac{1}{1+\lambda r_x(D)^2},\\
    \omega_{\mathrm{exp}}(D)
    &=
    \exp\{-\lambda r_x(D)^2\}.
\end{aligned}
\label{eq:controlled_taper_rules}
\end{equation}
All weighting values are computed from detached remoteness statistics;
the policy cannot reduce its loss by differentiating through the taper.

\paragraph{Near-field-retention matching.}
Let $\mathcal N$ denote the prespecified useful near-field set and define
its retained first-order update magnitude as
\begin{equation}
    \rho_{\mathrm{near}}(\omega)
    =
    \frac{
        \sum_{z\in\mathcal N}
        \omega(D_\theta(z))
        I